\documentclass[11pt,a4paper]{article}

% 语言和编码
\usepackage[utf8]{inputenc}
\usepackage[T1]{fontenc}
\usepackage[english]{babel}

% 数学公式
\usepackage{amsmath,amssymb,amsthm}
\usepackage{mathtools}

% 图表
\usepackage{graphicx}
\usepackage{booktabs}
\usepackage{multirow}
\usepackage{subcaption}

% 代码
\usepackage{listings}
\usepackage{xcolor}

% 页面布局
\usepackage[margin=2.5cm]{geometry}

% 超链接
\usepackage{hyperref}
\hypersetup{
    colorlinks=true,
    linkcolor=blue,
    citecolor=blue,
    urlcolor=blue
}

% 代码样式
\lstset{
    basicstyle=\ttfamily\small,
    keywordstyle=\color{blue},
    commentstyle=\color{green!60!black},
    stringstyle=\color{red},
    numbers=left,
    numberstyle=\tiny,
    numbersep=5pt,
    frame=single,
    breaklines=true,
    showstringspaces=false,
    tabsize=4
}

% 定理环境
\newtheorem{theorem}{Theorem}[section]
\newtheorem{lemma}[theorem]{Lemma}
\newtheorem{proposition}[theorem]{Proposition}
\newtheorem{corollary}[theorem]{Corollary}
\newtheorem{definition}[theorem]{Definition}

% 标题信息
\title{Algorithm Unrolling for Continuous Optimization Problems\\
\large Deep Learning Course Project}

\author{Student Name\\
\texttt{student@email.com}}

\date{\today}

\begin{document}

\maketitle

\begin{abstract}
This project explores \textbf{Algorithm Unrolling} techniques, transforming classical iterative optimization algorithms into learnable neural networks. We systematically implement three unrolled architectures: LISTA for LASSO sparse coding, ADMM-Net for low-rank matrix recovery, and PGD-Net for quadratic programming.

\textbf{Core contributions}:
\begin{enumerate}
    \item Propose \textbf{LISTA-Momentum} with learnable momentum mechanism to improve convergence
    \item Design \textbf{dimension-agnostic LISTA} that can handle arbitrary input dimensions
    \item Provide \textbf{spectral radius stability verification} and \textbf{parameter matrix analysis}
\end{enumerate}

\textbf{Results}: LISTA-Momentum outperforms classical ISTA on both in-distribution (ID) and out-of-distribution (OOD) tests, with only 1.06× generalization degradation. Dimension-agnostic LISTA achieves 3.77-8.72× performance improvement within training dimension range.

\bigskip
\noindent\textbf{Keywords:} Algorithm Unrolling, LISTA, Optimization, Deep Learning, Generalization
\end{abstract}

\tableofcontents
\newpage

% ==================== 引言 ====================
\section{Introduction}
\label{sec:intro}

\subsection{Background}

Iterative optimization algorithms are fundamental tools in machine learning and signal processing. However, traditional algorithms often require numerous iterations to converge, leading to high computational costs. \textbf{Algorithm Unrolling} is a technique that transforms iterative algorithms into deep networks by learning algorithm parameters to accelerate convergence \cite{gregor2010learning, monga2021algorithm}.

\subsection{Research Objectives}

This project systematically implements three unrolled architectures:
\begin{enumerate}
    \item \textbf{LASSO} — LISTA and its variants
    \item \textbf{Low-rank Matrix Recovery} — ADMM-Net and its variants
    \item \textbf{Quadratic Programming} — PGD-Net and its variants
\end{enumerate}

We focus on the following questions:
\begin{itemize}
    \item \textbf{Generalization}: Can unrolled networks generalize to unseen problem instances?
    \item \textbf{Dimension-agnostic}: Can we design a network that handles different input dimensions?
    \item \textbf{Theoretical analysis}: What about stability and convergence of unrolled networks?
\end{itemize}

\subsection{Main Contributions}

\begin{enumerate}
    \item \textbf{LISTA-Momentum}: Propose LISTA with learnable momentum to improve convergence and solve generalization
    \item \textbf{Dimension-agnostic LISTA}: Design dimension-agnostic architecture for arbitrary input dimensions
    \item \textbf{Theoretical analysis}: Provide spectral radius stability verification and parameter matrix analysis
    \item \textbf{Systematic experiments}: Conduct comprehensive comparison and ablation studies
\end{enumerate}

% ==================== 预备知识 ====================
\section{Preliminaries}
\label{sec:prelim}

\subsection{Proximal Operator}

The proximal operator is a key tool for solving non-smooth optimization problems \cite{parikh2014proximal}:
\begin{equation}
    \text{prox}_{\lambda f}(v) = \arg\min_x \left( f(x) + \frac{1}{2\lambda} \|x - v\|^2 \right)
\end{equation}

\subsection{Projection Operator}

Projection onto a convex set $C$:
\begin{equation}
    \text{Proj}_C(x) = \arg\min_{y \in C} \|y - x\|^2
\end{equation}

\subsection{ADMM Algorithm}

The Alternating Direction Method of Multipliers (ADMM) solves \cite{boyd2011distributed}:
\begin{equation}
    \min_{x,z} f(x) + g(z) \quad \text{s.t.} \quad Ax + Bz = c
\end{equation}

% ==================== 展开方法 ====================
\section{Unrolling Methods}
\label{sec:methods}

\subsection{LASSO — LISTA and Variants}

\textbf{Optimization problem}:
\begin{equation}
    \min_x \frac{1}{2} \|Ax - b\|^2 + \lambda \|x\|_1
\end{equation}

\textbf{ISTA iteration} \cite{beck2009fast}:
\begin{equation}
    x_{k+1} = \text{SoftThreshold}(x_k - \eta A^T(Ax_k - b), \eta\lambda)
\end{equation}

\subsubsection{LISTA (Gregor \& LeCun, 2010)}

Unroll $T$ ISTA iterations into $T$ layers with learnable parameters $W_1, W_2, \theta$:
\begin{equation}
    x_{k+1} = \sigma(W_1 b + W_2 x_k; \theta_k)
\end{equation}

\subsubsection{LISTA-CP (Chen et al., 2018)}

\textbf{Core improvement}: Couple weights using $B$:
\begin{equation}
    W_1 = \eta B, \quad W_2 = I - \eta BA
\end{equation}

\textbf{Problem}: $B$ is tied to training matrix $A_{train}$, cannot generalize to new $A$.

\subsubsection{LISTA-Momentum (Ours)}

\textbf{Core innovation}: Introduce learnable momentum:
\begin{equation}
\boxed{
\begin{aligned}
y_k &= x_k + \sigma(\beta) \cdot (x_k - x_{k-1}) & \text{(momentum)} \\
g_k &= A^T(A y_k - b) & \text{(gradient with current A)} \\
x_{k+1} &= \text{SoftThreshold}(y_k - \eta \cdot W g_k; \theta) & \text{(learnable transform)}
\end{aligned}
}
\end{equation}

\textbf{Key design}:
\begin{enumerate}
    \item Don't embed A in weights: compute gradient with current A at each layer
    \item Learn general gradient transformation rule W: independent of A
    \item Learnable momentum parameter $\beta$: adaptive acceleration
\end{enumerate}

\subsubsection{Dimension-agnostic LISTA (Ours)}

\textbf{Core innovation}: Design dimension-agnostic architecture:
\begin{enumerate}
    \item LayerNorm normalization: eliminate dimension effects
    \item Shared MLP transformation: same transform for each element
    \item Scale factor: restore pre-normalization scale
\end{enumerate}

\subsection{Low-rank Matrix Recovery — ADMM-Net}

\textbf{Optimization problem}:
\begin{equation}
    \min_X \|X\|_* \quad \text{s.t.} \quad P_\Omega(X) = P_\Omega(M)
\end{equation}

\textbf{ADMM-Net unrolling} \cite{sun2016deep}:
\begin{itemize}
    \item ADMMNet: Basic ADMM unrolling
    \item ADMMNetV2: With learnable linear transforms
    \item SoftImputeNet: Soft-Impute unrolling
\end{itemize}

\subsection{Quadratic Programming — PGD-Net}

\textbf{Optimization problem}:
\begin{equation}
    \min_x \frac{1}{2} x^T Q x + c^T x \quad \text{s.t.} \quad x \in C
\end{equation}

\textbf{PGD-Net variants}:
\begin{itemize}
    \item PGDNet: Basic PGD unrolling
    \item PGDMomentumNet: With Nesterov momentum
    \item PGDLearnedQNet: Learning Q matrix correction
\end{itemize}

% ==================== 实验设置 ====================
\section{Experimental Setup}
\label{sec:setup}

\subsection{Data Generation}

\textbf{LASSO}:
\begin{itemize}
    \item Measurement matrix $A \in \mathbb{R}^{m \times n}$ with normalized columns
    \item Sparse signal $x$ with sparsity $s = 10$
    \item Observation $b = Ax + \epsilon$, $\epsilon \sim \mathcal{N}(0, 0.001^2)$
\end{itemize}

\subsection{Training Configuration}

\begin{table}[h]
\centering
\begin{tabular}{lc}
\toprule
\textbf{Parameter} & \textbf{Value} \\
\midrule
Optimizer & Adam \\
Learning rate & 1e-3 \\
LR scheduler & ReduceLROnPlateau \\
Gradient clipping & max\_norm=1.0 \\
Early stopping & patience=20 \\
Batch size & 64 \\
Training samples & 1000 \\
Validation samples & 200 \\
Test samples & 100 \\
\bottomrule
\end{tabular}
\caption{Training configuration}
\label{tab:config}
\end{table}

\subsection{Evaluation Metrics}

\begin{itemize}
    \item \textbf{Relative error}: $\|x_{\text{true}} - x_{\text{pred}}\| / \|x_{\text{true}}\|$
    \item \textbf{Support recovery rate} (LASSO)
    \item \textbf{Rank recovery} (low-rank matrix)
    \item \textbf{Constraint violation} (QP)
    \item \textbf{Optimality gap}: $(f(x) - f^*) / |f^*|$
    \item \textbf{Inference time}: Wall-clock time (ms)
    \item \textbf{Spectral radius}: $\rho(W_2) = \max_i |\lambda_i(I - \eta BA)|$
\end{itemize}

% ==================== 实验结果 ====================
\section{Experimental Results}
\label{sec:results}

\subsection{LASSO Experiment}

\subsubsection{Fair Comparison (Same Iterations)}

\begin{table}[h]
\centering
\begin{tabular}{ccccc}
\toprule
\textbf{T} & \textbf{ISTA} & \textbf{FISTA} & \textbf{LISTA-CP} & \textbf{LISTA-Momentum} \\
\midrule
5 & 0.859 & 0.852 & 0.517 & \textbf{0.643} \\
10 & 0.843 & 0.821 & 0.448 & \textbf{0.526} \\
20 & 0.821 & 0.742 & 0.360 & \textbf{0.447} \\
\bottomrule
\end{tabular}
\caption{Fair comparison at same number of iterations}
\label{tab:fair}
\end{table}

\subsubsection{Generalization Comparison}

\begin{table}[h]
\centering
\begin{tabular}{lcccc}
\toprule
\textbf{Method} & \textbf{ID} & \textbf{OOD} & \textbf{Degradation} & \textbf{vs ISTA} \\
\midrule
ISTA & 0.843 & 0.844 & 1.00× & Baseline \\
FISTA & 0.821 & 0.822 & 1.00× & 1.03× \\
LISTA-CP & 0.448 & 1.215 & 2.71× & 1.88× (ID) \\
\textbf{LISTA-Momentum} & \textbf{0.526} & \textbf{0.555} & \textbf{1.06×} & \textbf{1.60×} \\
\bottomrule
\end{tabular}
\caption{Generalization comparison}
\label{tab:gen}
\end{table}

\subsection{Dimension Generalization}

\begin{table}[h]
\centering
\begin{tabular}{lcccc}
\toprule
\textbf{Dimension n} & \textbf{ISTA} & \textbf{LISTA-DimAgnostic} & \textbf{vs ISTA} & \textbf{Status} \\
\midrule
50 & 0.479 & \textbf{0.088} & \textbf{5.44×} & Training \\
100 & 0.701 & \textbf{0.080} & \textbf{8.72×} & Training, Best \\
150 & 0.799 & \textbf{0.106} & \textbf{7.52×} & Training \\
200 & 0.852 & \textbf{0.158} & \textbf{5.38×} & Training \\
300 & 0.900 & \textbf{0.239} & \textbf{3.77×} & Training \\
\bottomrule
\end{tabular}
\caption{Dimension generalization results}
\label{tab:dim}
\end{table}

\subsection{Parameter Matrix W Analysis}

\begin{table}[h]
\centering
\begin{tabular}{ccccc}
\toprule
\textbf{Layer} & \textbf{Cosine Sim.} & \textbf{Spectral Radius} & \textbf{Condition Num.} & \textbf{Effective Rank} \\
\midrule
1 & 0.915 & 1.452 & 3.01 & 200 \\
5 & 0.929 & 1.522 & 2.23 & 200 \\
10 & 0.964 & 1.426 & 2.29 & 200 \\
\bottomrule
\end{tabular}
\caption{Parameter matrix W analysis}
\label{tab:w_analysis}
\end{table}

% ==================== 消融实验 ====================
\section{Ablation Studies}
\label{sec:ablation}

\subsection{Effect of Layer Count}

\begin{table}[h]
\centering
\begin{tabular}{ccccc}
\toprule
\textbf{T} & \textbf{ISTA} & \textbf{LISTA-CP} & \textbf{LISTA-Momentum} & \textbf{vs ISTA} \\
\midrule
5 & 0.859 & 0.517 & 0.643 & 1.34× \\
10 & 0.843 & 0.448 & 0.526 & 1.60× \\
20 & 0.821 & 0.360 & 0.447 & 1.84× \\
\bottomrule
\end{tabular}
\caption{Effect of number of layers}
\label{tab:layers}
\end{table}

\subsection{Noise Robustness}

\begin{table}[h]
\centering
\begin{tabular}{ccc}
\toprule
\textbf{Noise Level $\sigma$} & \textbf{LISTA-Momentum} & \textbf{vs ISTA} \\
\midrule
0.001 & 0.542 & 1.55× \\
0.005 & 0.543 & 1.55× \\
0.010 & 0.543 & 1.55× \\
0.020 & 0.546 & 1.54× \\
0.050 & 0.562 & 1.50× \\
0.100 & 0.610 & 1.38× \\
\bottomrule
\end{tabular}
\caption{Noise robustness}
\label{tab:noise}
\end{table}

% ==================== 结论 ====================
\section{Conclusion}
\label{sec:conclusion}

\subsection{Main Contributions}

\begin{enumerate}
    \item \textbf{LISTA-Momentum}: Propose LISTA with learnable momentum to improve convergence and solve generalization
    \item \textbf{Dimension-agnostic LISTA}: Design dimension-agnostic architecture for arbitrary input dimensions
    \item \textbf{Theoretical analysis}: Provide spectral radius stability verification and parameter matrix analysis
    \item \textbf{Systematic experiments}: Conduct comprehensive comparison and ablation studies
\end{enumerate}

\subsection{Key Findings}

\begin{enumerate}
    \item \textbf{LISTA-Momentum generalization}: By not embedding A in weights, achieve only 1.06× generalization degradation
    \item \textbf{Dimension-agnostic success}: Using LayerNorm + shared MLP, achieve 3.77-8.72× performance improvement
    \item \textbf{Parameter matrix W}: W is close to identity matrix, learning ``how to improve gradient''
\end{enumerate}

\subsection{Future Work}

\begin{enumerate}
    \item \textbf{Extended verification}: Verify on larger-scale problems
    \item \textbf{Theoretical analysis}: Study convergence and generalization
    \item \textbf{Application推广}: Apply general architecture to ADMM-Net and PGD-Net
\end{enumerate}

% ==================== 参考文献 ====================
\begin{thebibliography}{10}

\bibitem{gregor2010learning}
Gregor, K., \& LeCun, Y. (2010). Learning fast approximations of sparse coding. In \textit{ICML}.

\bibitem{chen2018theoretical}
Chen, X., Liu, J., Wang, Z., \& Yin, W. (2018). Theoretical linear convergence of unfolded ISTA and its practical weights and thresholds. In \textit{NeurIPS}.

\bibitem{beck2009fast}
Beck, A., \& Teboulle, M. (2009). A fast iterative shrinkage-thresholding algorithm for linear inverse problems. \textit{SIAM Journal on Imaging Sciences}, 2(1), 183-202.

\bibitem{boyd2011distributed}
Boyd, S., Parikh, N., Chu, E., Peleato, B., \& Eckstein, J. (2011). Distributed optimization and statistical learning via the alternating direction method of multipliers. \textit{Foundations and Trends in Machine Learning}, 3(1), 1-122.

\bibitem{parikh2014proximal}
Parikh, N., \& Boyd, S. (2014). Proximal algorithms. \textit{Foundations and Trends in Optimization}, 1(3), 127-239.

\bibitem{duchi2008efficient}
Duchi, J., Shalev-Shwartz, S., Singer, Y., \& Chandra, T. (2008). Efficient projections onto the ℓ1-ball for learning in high dimensions. In \textit{ICML}.

\bibitem{monga2021algorithm}
Monga, V., Li, Y., \& Eldar, Y. C. (2021). Algorithm unrolling: Interpretable, efficient deep learning for signal and image processing. \textit{IEEE Signal Processing Magazine}, 38(2), 18-44.

\bibitem{sun2016deep}
Sun, J., Li, H., \& Xu, Z. (2016). Deep ADMM-Net for compressive sensing MRI. In \textit{NeurIPS}.

\bibitem{andrychowicz2016learning}
Andrychowicz, M., Denil, M., Gomez, S., Hoffman, M. W., Pfau, D., Schaul, T., \& de Freitas, N. (2016). Learning to learn by gradient descent by gradient descent. In \textit{NeurIPS}.

\bibitem{borgerding2017amp}
Borgerding, M., Schniter, P., \& Rangan, S. (2017). AMP-inspired deep networks for sparse linear inverse problems. \textit{IEEE Transactions on Signal Processing}, 65(18), 4293-4308.

\end{thebibliography}

\end{document}
